% from ega4-11-fr.tex
\oldpage[IV-3]{171}
\begin{definition}[11.10.2]
\label{IV.11.10.2-fr}
Lorsque les conditions équivalentes de (11.10.1) sont vérifiées, on dit que la
famille $(f_\lambda)$ est \emph{schématiquement dominante}. Lorsque les
$f_\lambda$ sont les immersions canoniques dans $X$ d'une famille
$(Z_\lambda)$ de sous-préschémas de $X$, on dit aussi que la famille
$(Z_\lambda)$ est \emph{schématiquement dense}.
\end{definition}

\begin{remarks}[11.10.3]
\label{IV.11.10.3-fr}
\begin{enumerate}[label=(\roman*)]
  \item La notion de famille schématiquement dominante est \emph{locale} sur
  $X$, comme il résulte par exemple de la forme b) dans (11.10.1)~: si
  $(W_\alpha)$ est un recouvrement ouvert de $X$, la famille $(f_\lambda)$ est
  schématiquement dominante si et seulement s'il en est ainsi de chacune des
  familles formées des morphismes
  $f_\lambda^{-1}(W_\alpha)\to W_\alpha$, restrictions des $f_\lambda$.
  \item Si $Z$ est le préschéma somme des $Z_\lambda$, $f:Z\to X$ le morphisme
  coïncidant avec $f_\lambda$ dans chacun des $Z_\lambda$, il revient au même
  de dire que la famille $(f_\lambda)$ est schématiquement dominante, ou que
  la famille réduite au seul élément $f$ l'est.
  \item Soit $M$ un second ensemble d'indices et, pour tout $\lambda\in L$,
  soit $(g_{\lambda\mu})_{\mu\in M}$ une famille de morphismes
  $g_{\lambda\mu}:Z'_{\lambda\mu}\to Z_\lambda$~; si, pour chaque
  $\lambda\in L$, la famille $(g_{\lambda\mu})_{\mu\in M}$ est
  schématiquement dominante, alors, il revient au même de dire que la famille
  $(f_\lambda)_{\lambda\in L}$ est schématiquement dominante, ou que la
  famille
  $(f_\lambda\circ g_{\lambda\mu})_{(\lambda,\mu)\in L\times M}$ l'est
  (11.9.2).
  \item Soit $f:Z\to X$ un morphisme tel que $f_*(\mathcal O_Z)$ soit un
  $\mathcal O_X$-Module quasi-cohérent (par exemple un morphisme quasi-compact
  et quasi-séparé (1.7.4)). Alors, dire que $f$ est schématiquement dominant
  signifie que \emph{l'image fermée de $Z$ par $f$} ($\mathbf I$, 9.5.3) est
  \emph{identique à $X$}.
\end{enumerate}
\end{remarks}

\begin{proposition}[11.10.4]
\label{IV.11.10.4-fr}
Si la famille de morphismes $f_\lambda:Z_\lambda\to X$ est schématiquement
dominante, la réunion des $f_\lambda(Z_\lambda)$ est dense dans $X$.
Inversement, si cette réunion est dense dans $X$ et si de plus $X$ est réduit,
la famille $(f_\lambda)$ est schématiquement dominante.
\end{proposition}

\begin{proof}
La première assertion résulte aussitôt de (11.10.1, b)). D'autre part, si $X$
est réduit, et si la réunion des $f_\lambda(Z_\lambda)$ est dense dans $X$,
alors, si l'on a des factorisations (11.10.1.2) pour tout $\lambda\in L$, on a
aussi des factorisations
\[
(f_\lambda^{-1}(U))_{\mathrm{red}}
\xrightarrow{(g_\lambda)_{\mathrm{red}}}Y_{\mathrm{red}}
\xrightarrow{j_{\mathrm{red}}}U,
\]
et l'hypothèse entraîne que l'espace sous-jacent à $Y$ est identique à $U$,
donc $Y=Y_{\mathrm{red}}=U$ puisque $U$ est réduit.
\end{proof}

Les résultats du n$^\circ$~11.9 sur les familles séparantes se traduisent en
résultats sur les familles schématiquement dominantes~:

\begin{theorem}[11.10.5]
\label{IV.11.10.5-fr}
Soient $(f_\lambda)$ une famille de morphismes
$f_\lambda:Z_\lambda\to X$, $g:X'\to X$ un morphisme, et posons
$Z'_\lambda=Z_\lambda\times_XX'$,
$f'_\lambda=(f_\lambda)_{(X')}:Z'_\lambda\to X'$.
\begin{enumerate}[label=(\roman*)]
  \item Si $g$ est fidèlement plat et si la famille $(f'_\lambda)$ est
  schématiquement dominante, alors la famille $(f_\lambda)$ est
  schématiquement dominante.
  \item Supposons que $g$ soit un morphisme plat et en outre que l'une des
  conditions suivantes soit vérifiée~:
  \begin{enumerate}[label=\emph{\alph*)}]
    \item $L$ est fini et les $f_\lambda$ sont quasi-compacts.
    \item Le morphisme $g$ est localement de présentation finie.
  \end{enumerate}
  Alors, si la famille $(f_\lambda)$ est schématiquement dominante, il en est
  de même de la famille $(f'_\lambda)$.
\end{enumerate}
\end{theorem}

